% Options for packages loaded elsewhere
\PassOptionsToPackage{unicode}{hyperref}
\PassOptionsToPackage{hyphens}{url}
\documentclass[
]{article}
\usepackage{xcolor}
\usepackage[margin=1in]{geometry}
\usepackage{amsmath,amssymb}
\setcounter{secnumdepth}{-\maxdimen} % remove section numbering
\usepackage{iftex}
\ifPDFTeX
  \usepackage[T1]{fontenc}
  \usepackage[utf8]{inputenc}
  \usepackage{textcomp} % provide euro and other symbols
\else % if luatex or xetex
  \usepackage{unicode-math} % this also loads fontspec
  \defaultfontfeatures{Scale=MatchLowercase}
  \defaultfontfeatures[\rmfamily]{Ligatures=TeX,Scale=1}
\fi
\usepackage{lmodern}
\ifPDFTeX\else
  % xetex/luatex font selection
\fi
% Use upquote if available, for straight quotes in verbatim environments
\IfFileExists{upquote.sty}{\usepackage{upquote}}{}
\IfFileExists{microtype.sty}{% use microtype if available
  \usepackage[]{microtype}
  \UseMicrotypeSet[protrusion]{basicmath} % disable protrusion for tt fonts
}{}
\makeatletter
\@ifundefined{KOMAClassName}{% if non-KOMA class
  \IfFileExists{parskip.sty}{%
    \usepackage{parskip}
  }{% else
    \setlength{\parindent}{0pt}
    \setlength{\parskip}{6pt plus 2pt minus 1pt}}
}{% if KOMA class
  \KOMAoptions{parskip=half}}
\makeatother
\usepackage{graphicx}
\makeatletter
\newsavebox\pandoc@box
\newcommand*\pandocbounded[1]{% scales image to fit in text height/width
  \sbox\pandoc@box{#1}%
  \Gscale@div\@tempa{\textheight}{\dimexpr\ht\pandoc@box+\dp\pandoc@box\relax}%
  \Gscale@div\@tempb{\linewidth}{\wd\pandoc@box}%
  \ifdim\@tempb\p@<\@tempa\p@\let\@tempa\@tempb\fi% select the smaller of both
  \ifdim\@tempa\p@<\p@\scalebox{\@tempa}{\usebox\pandoc@box}%
  \else\usebox{\pandoc@box}%
  \fi%
}
% Set default figure placement to htbp
\def\fps@figure{htbp}
\makeatother
\setlength{\emergencystretch}{3em} % prevent overfull lines
\providecommand{\tightlist}{%
  \setlength{\itemsep}{0pt}\setlength{\parskip}{0pt}}
\usepackage{bookmark}
\IfFileExists{xurl.sty}{\usepackage{xurl}}{} % add URL line breaks if available
\urlstyle{same}
\hypersetup{
  hidelinks,
  pdfcreator={LaTeX via pandoc}}

\author{}
\date{\vspace{-2.5em}}

\begin{document}

\section{Materials and methods}\label{materials-and-methods}

\subsection{Study sites}\label{study-sites}

The experiment was conducted in the Phoenix metropolitan area of central
Arizona, USA (33.5° N, 111.8° W). The climate of the region is hot and
arid with a mean annual temperature of 22.5°C and total precipitation of
19.5 cm (Cervany 1996). Visible and UV irradiance are relatively high in
central Arizona because of its low latitude and the prevalence of clear,
dry skies for much of the year. For example, Phoenix receives on average
85\% of the annual total possible sunshine (Cervany 1996).

The first study site was an open, unshaded, second-story balcony
(\(3.6  \text{m} \times 2.0  \text{m}\)) on the southeast corner of a
house in Scottsdale, AZ, a suburb of Phoenix, referred to as the
`balcony' site. In order to better mimic the natural conditions under
which \emph{L. tridentata} litter decomposes, we repeated the experiment
at a second site in native \emph{L. tridentata} desert at the Desert
Botanical Gardens, Phoenix, AZ, 21 km from the first site. We refer to
this as the `desert' site. This site is in a native desert conservation
area and consists of scattered \emph{L. tridentata} shrubs (0.05
individuals/m2, 18\% foliar cover) with bare ground between shrubs (Day
et al.~2002). Our rationale for conducting the first experiment at the
balcony site was: (1) This site was easily accessed on a daily basis
which allowed us to work out treatment protocols in preliminary trials.
An optimal design for litterbags was developed during initial trials
here. (2) We were concerned about dust/soil contamination that we
suspected could make subsequent chemical analyses of litter difficult.
The balcony site provided a cleaner environment for litter where
contamination might not be as pronounced as on the desert floor. (3) A
comparison of results from the balcony with the desert site would
provide an indirect means to assess how important contact with the soil
surface (and subsequent inoculation with microbes) is in litter
decomposition in our system.

\subsection{Litter source and
preparation}\label{litter-source-and-preparation}

Recently abscised current-year \emph{L. tridentata} litter was collected
in November 2002 from beneath the canopy of several \emph{L. tridentata}
bushes growing in native desert adjacent to Arizona State University's
main campus in Tempe, AZ (within 25 km of both study sites). Rocks and
other non \emph{L. tridentata} material, as well as large twigs
(\textgreater2 mm in diameter), were removed through sieving and hand
sorting. The litter was furthe sorted into three litter types: (1) an
unaltered mixture of the leaves, twigs, and seeds, (2) twigs (\textless2
mm diameter) only, and (3) leaves only. Sorted litter was air dried and
stored in sealed plastic bags in the laboratory until placement in
litterbags. Just prior to placement in litterbags, litter was air dried
again for 2 weeks, oven dried at 30°C for 3 days, and then weighed and
placed in litterbags.

\subsection{Litterbags}\label{litterbags}

Litter was placed in \(10 \times 10 \text{cm}\) bags or envelopes
constructed of either 125 \unit{\micro\metre}-thick UVB-transparent film
(Aclar Type 22A film, Proplastics, Linden, NJ; transmission
\textgreater90\% through the UV-B and UV-A (315--400 nm) wavebands) or
125 \unit{\micro\metre}-thick UV-B-opaque film (Mylar-type Cadco clear
polyester film, Cadillac Plastic and Chemical, Phoenix, AZ; sharp
transmission cut-off below 325 nm). The Aclar envelopes represented a
`near-ambient UV-B' treatment while the Mylar-type envelopes represented
a `reduced UV-B' treatment. To our knowledge, neither of these film
materials have the unintended and potentially toxic effect that
cellulose acetate film has on some organisms (Krizek and Mirecki 2004).
Each envelope was comprised of a top and bottom piece of the appropriate
film, with the edges of the envelopes sealed with UV-B transparent tape
(Scotch Multitask tape, transmission \textgreater90\% through the UVB
and UVA wavebands, 3M, St.~Paul, MN). To allow water and microbes to
reach the litter, we drilled 2-mm diameter holes, spaced 1 cm apart,
throughout the top and bottom pieces of the film. We measured UV-B
irradiance directly under the top film of envelopes under clear skies at
midday in January with a UV-scanning spectroradiometer (OL 754, Optronic
Laboratories, Orlando, FL). Biologically effective solar UV-B dose
(based on the generalized plant damage action spectrum normalized to 300
nm, Caldwell 1971), was 85\% of ambient in Aclar envelopes (near-ambient
treatment) and 15\% of ambient in Mylar-type envelopes (reduced
treatment). To prevent smaller pieces of litter from falling through the
holes on the bottom side of the envelope, a
\(10 \text{cm} \times 10 \text{cm}\) piece of white nylon fabric
(athletic jersey mesh) was taped on the inside bottom. We placed 1.80
(±0.05) g of oven dried litter of a particular type in each envelope on
top of the nylon fabric.

\subsection{Experiments}\label{experiments}

At the balcony site, envelopes were placed in square, black plastic
greenhouse trays
(\(42 \text{cm} \times 43.5 \text{cm} \times 5.5 \text{cm}\) high
sidewalls) having a lattice bottom that allowed water to drain freely.
The trays were anchored on top of wooden frames (8.9 cm high), which
secured the trays to the balcony floor but allowed water drainage. Clear
monofilament fishing line spaced 5 cm apart was threaded across the
trays to secure litter envelopes in the bottom of the plastic trays. A
total of 60 envelopes were randomly assigned to seven trays on the
balcony. This represented a \(3 \times 2\) factorial design consisting
of three litter types (full mixture, twigs, or leaves) and two UV- B
levels (reduced or near-ambient) for each litter type. There were 10
replicate envelopes of each treatment/litter type combination. Envelopes
were placed at this site on 16 September, 2003 and retrieved on 22
February, 2004 after 21 weeks. Average air temperature and total
precipitation over this period were 18.9°C and 5.11 cm, respectively,
based on National Weather Service measurements at nearby Sky Harbor
International Airport.

At the desert site, envelopes were placed directly on the ground in a
level interstice between \emph{L. tridentata} shrubs so as to receive
full sun for most of the daytime. The interstice was divided into a grid
and envelopes were randomly assigned to grid locations. Envelopes were 8
cm apart, and adjacent shrubs were at least 1.5 m from the edges of the
grid. The envelopes at the desert site were modified by adding 2-cm long
flaps of Aclar that extended from each corner. Nails (13 cm long) were
driven through the flaps into the soil to secure the envelopes to the
ground. As in the first experiment, 60 envelopes were used, representing
10 replicates of each treatment/litter type combination. Envelopes were
placed in the desert on 22 December, 2003 and retrieved on 16 April,
2004 after 16 weeks. Average air temperature and total precipitation
over this period were 17.2°C and 10.44 cm, respectively.

After collection, litter was sieved to remove inorganic material, air
dried for 2 weeks, oven dried at 30°C for 3 days and weighed. We
measured the total ash content of litter samples by placing subsamples
collected before and after field treatments in a muffle furnace at 550°C
for 8 h. The ash proportion of samples was used to correct mass values
for soil and dust contamination. Decomposition rate was expressed as
mass loss and expressed as percent loss over the treatment period.

\subsection{Chemical composition}\label{chemical-composition}

We assessed the chemical composition of the leaf litter at the balcony
site prior to placing litterbags outdoors and at the end of the
treatment period. Specifically, we measured concentrations of
holocellulose, lignin, soluble fats and lipids, and total organic C and
N. We did not assess the composition of the other litter types from the
balcony site. We attempted to assess the composition of litter from the
desert site, but suspect that slight contamination with dust prevented
us from obtaining reliable results. Hence, only chemical composition of
the balcony leaf litter is reported.

We modified the techniques of Allen (1989) for isolation and
quantification of fats and lipids, lignin and holocellulose. Dried
litter was ground to a fine powder using a ball-mill grinder and weighed
with an analytical balance and samples were further divided for analysis
of fats and lipids, lignin, holocellulose, total organic C and N, and
ash. For fats and lipids, and lignin, litter samples ( 0.75 g) were
placed on glass--fiber filter paper and tied into bundles using nylon
thread, and extracted in diethyl ether for 6 h using a Soxhlet
apparatus. Samples were dried (30°C) overnight, cooled and weighed for
fat and lipid content. Soluble carbohydrates were removed by boiling the
samples in water for 3 h. Proteins were removed by adding 16.5 ml of
10\% H2SO4 and boiling for an additional hour. Samples were allowed to
settle and the supernatant was removed. Hydrolysis of lignin was
accomplished by adding 12 ml of 72\% H2SO4 for 2 h with occasional
stirring. Acid strength was reduced to 3\% by the addition of water, and
the samples were boiled for another 4 h. Samples were first filtered and
washed repeatedly with hot water, then dried (30°C) overnight, cooled
and weighed on an analytical balance for lignin determination.

For holocellulose, samples ( 0.25 g) were placed in 50-ml Erlenmeyer
flasks containing 15 ml of water. Samples were delignified by the
addition of 0.15 g of sodium chlorite (NaClO2) and 0.5 ml of 10\% acetic
acid. Flasks were covered with a glass marble and heated at 75°C for 4
h. Additional NaClO2 and acetic acid were added at hourly intervals
during heating. Samples were filtered and washed repeatedly with chilled
water, acetone, and diethyl ether. Samples were then dried (30°)
overnight, cooled, and weighed on an analytical balance. Concentrations
were corrected for total ash content by placing subsamples in a muffle
furnace at 550°C for 8 h. Total organic C and N were measured with a
dynamic flash combustion analyzer (FlashEA 1112, Thermo Finnigan,
Waltham, MA).

\subsection{Statistical analyses}\label{statistical-analyses}

We examined the effect of UV-B treatment on mass loss and chemical
composition of litter using ANOVA, followed by comparison of initial,
and reduced and near-ambient treatment means using the least significant
difference test. All data were transformed using an arc-sine
transformation to meet the assumption of homogeneity of variances.

\end{document}
